\documentclass[12pt]{article}

\setlength{\parindent}{0pt}
\setcounter{secnumdepth}{3}
\usepackage[margin=1in]{geometry}

% packages
\usepackage{amsmath, amsfonts, amsthm, amssymb}
\usepackage{graphicx, float}
\usepackage{mathtools}
\usepackage{titlesec}
\usepackage{interval}
\usepackage{hyperref}
\usepackage{siunitx}
\usepackage{titling}
\usepackage{vwcol}
\usepackage{setspace}
\usepackage{empheq}
\usepackage{cancel}
\usepackage{esdiff}
\usepackage{multicol}
\usepackage{mdframed}
\usepackage{esdiff}
\usepackage{tikzsymbols}
\usepackage{multicol}
\usepackage{tikz}
\usepackage{varwidth}
\usepackage{pgfplots}
\usepackage{fancyhdr}
\usepackage{enumitem}
\usepackage{tcolorbox}

\intervalconfig {
	soft open fences
}

% header
\pagestyle{fancy}
\fancyhf{}
\lhead{Ethan Fan}
\rhead{AP Calculus BC}
\cfoot{\thepage}

% section format
\titleformat{\section}
  {\large \bfseries}
  {}
  {0em}
  {}
\titlespacing*{\section}{0pt}{0.5em}{0.3em}

\titleformat{\subsection}[runin]
  {\bfseries}
  {}
  {0em}
  {}
\titlespacing*{\subsection}{0pt}{0pt}{0em}

\titleformat{\subsubsection}[runin]
  {\itshape}
  {(\roman{subsubsection})}
  {0em}
  {}

% commands
\newcommand{\alignedintertext}[1]{%
  \noalign{%
    \vskip\belowdisplayshortskip
    \vtop{\hsize=\linewidth#1\par
    \expandafter}%
    \expandafter\prevdepth\the\prevdepth
  }%
}

\newcommand*{\paren}[1]{\ensuremath\left(#1\right)}
\newcommand*{\limit}[2][x]{\ensuremath{\displaystyle\lim_{#1\to#2}}}
\newcommand*{\Limit}[3][x]{\ensuremath{\displaystyle\lim_{#1\to#2}\left[#3\right]}}
\newcommand*{\deriv}[1][x]{\ensuremath{\dfrac{\mathrm{d}}{\mathrm{d}#1}}}
\newcommand*{\Deriv}[2][x]{\ensuremath{\dfrac{\mathrm{d}}{\mathrm{d}#1}\left[#2\right]}}
\newcommand{\ddx}{\dfrac{d}{dx}}
\newcommand{\dydx}{\dfrac{dy}{dx}}
\newcommand{\dydt}{\dfrac{dy}{dt}}
\newcommand{\dxdt}{\dfrac{dx}{dt}}
\newcommand*{\iinteg}[2][x]{\ensuremath{\displaystyle\int #2\;\mathrm{d}#1}}
\newcommand*{\dinteg}[4][x]{\ensuremath{\displaystyle\int_{#2}^{#3}#4\;\mathrm{d}#1}}
\newcommand*{\abs}[1]{\ensuremath{\left|#1\right|}}
\newcommand*{\eps}{\varepsilon}
\newcommand*{\real}{\ensuremath\mathbb{R}}
\newcommand*{\integer}{\ensuremath\mathbb{Z}}
\newcommand*{\posint}{\ensuremath\mathbb{Z^+}}
\newcommand*{\cbrt}[1]{\ensuremath{\sqrt[3]{#1}}}
\newcommand*{\lheqzero}{\ensuremath{\underset{\text{L'H}}{\overset{\left[\frac00\right]}{=}}}}
\newcommand*{\lheqinfty}{\ensuremath{\underset{\text{L'H}}{\overset{\left[\frac{\infty}{\infty}\right]}{=}}}}
\newcommand{\tor}{\text{ or }}
\newcommand{\floor}[1]{\lfloor #1 \rfloor}

\newcommand{\vem}{\vspace{0.5em}}
\newcommand{\example}[1]{\section{Example #1}}
\newcommand{\problem}[1]{\section{Problem #1}}
\newcommand{\subproblem}[1]{\subsection{(#1)} \,}

\DeclareMathOperator{\dne}{DNE}
\DeclareMathOperator{\sgn}{sgn}

\newcommand{\definition}[1]{\begin{tcolorbox}[colback=red!5!white,colframe=red!75!black,parbox=false] #1 \end{tcolorbox}}

\newcommand{\theorem}[2]{\begin{tcolorbox}[title={#1},colback=blue!5!white,colframe=blue!75!black,parbox=false] #2 \end{tcolorbox}}

\allowdisplaybreaks
% \postdisplaypenalty=100000

% document
\begin{document}

\vspace*{0cm}

% opening
\begin{center}
    {\Large \bfseries Problem Set 10} \\[0.5cm]
    {\large Differentials and Linearization} \\[0.35cm]
    {\normalsize September 28, 2026}
\end{center}

% problems

\problem{1}

\subproblem{d}
\begin{gather*}
    \lim_{dx\to 0} \frac{\Delta y-dy}{d^2x}=\lim_{dx\to 0} 7+dx=7\\
    f''(2)=6\cdot 2+2=14\\
    \lim_{dx\to 0} \frac{\Delta y-dy}{d^2x} =\frac{1}{2}f''(2)\\
    \Delta y-dy= \frac{1}{2}f''(2) d^2x\\
    |f(x)-L(x)|\le \frac{M}{2}(x-a)^2
\end{gather*}

\problem{2}

\subproblem{a}
\begin{gather*}
    y=x^2\sin 2x\\
    y'=2x\sin 2x+2x^2\cos 2x\\
    dy=\boxed{2x(\sin 2x+x\cos 2x)\, dx}
\end{gather*}

\subproblem{b}
\begin{gather*}
    y=\frac{u+1}{u-1}\\
    y'=\frac{-2}{(u-1)^2}\\
    dy=\boxed{-\frac{2}{(u-1)^2}\, du}
\end{gather*}

\subproblem{c}
\begin{gather*}
    y=e^{\tan \pi t}\\
    y'=\pi \sec^2 (\pi t)e^{\tan \pi t}\\
    dy=\boxed{\pi \sec^2 (\pi t)e^{\tan \pi t}\, dt}
\end{gather*}

\subproblem{d}
\begin{gather*}
    y=\ln (1+\sqrt x)\\
    y'=\frac{1}{1+\sqrt{x}}\cdot \frac{1}{2\sqrt{x}}\\
    dy=\boxed{\frac{1}{2\sqrt{x}+2x}\, dx}
\end{gather*}

\subproblem{e}
\begin{gather*}
    y=\arctan \frac{1}{x}\\
    y'=\frac{1}{1+\frac{1}{x^2}} \cdot \frac{-1}{x^2}\\
    dy=\boxed{-\frac{1}{x^2+1}\, dx}
\end{gather*}

\subproblem{f}
\begin{gather*}
    d(\arctan x)=\frac{1}{1+x^2}\, dx\\
    d(\arctan \frac{1}{x})+d(\arctan x)=0
\end{gather*}
$\arctan x+\arctan \frac{1}{x}$ is constant on both $(0,\infty)$ and $(-\infty ,0)$.\vem

\subproblem{g}
\begin{gather*}
    y=e^\frac{x}{10}\\
    y'=\frac{1}{10}e^\frac{x}{10} \implies y'(0)=\frac{1}{10}\\
    dy=\frac{1}{10}\cdot 0.1=0.01\\
    \Delta y = e^{0.01}-1\approx 0.010005\\
    |\Delta y - dy|=0.000005\\
    y=\tan x\\
    y'=\sec^2 x \implies y'(\frac{\pi}{4})=2\\
    dy=2\cdot -0.1 = -0.2\\
    \Delta y = \tan (\frac{\pi}{4}-0.1) - 1\approx -0.18237\\
    |\Delta y - dy|=0.01763
\end{gather*}
The first approximation is much better, as its second derivative's magnitude is closer to zero.\vem

\subproblem{h}
\begin{gather*}
    x^2+xy+y^2=7\\
    2x\, dx + x\, dy+y\, dx+2y\, dy=0\\
    dy=-\frac{2x+y}{2y+x}\, dx = - \frac{4}{5}\, dx\\
    y=-\frac{4}{5}(x-1)+2=1.92\\
    y=1.918
\end{gather*}

\problem{4}

\subproblem{a}
\begin{gather*}
    f(-1)=4\\
    f'(-1)=-4-6=-10\\
    y-4=-10(x+1)
\end{gather*}

\subproblem{b}
\begin{gather*}
    f(\frac{\pi}{2})=0\\
    f'(\frac{\pi}{2})=-1\\
    y=-(x-\frac{\pi}{2})
\end{gather*}

\subproblem{c}
\begin{gather*}
    f(1)=0\\
    f'(1)=1\\
    y=x-1
\end{gather*}

\subproblem{d}
\begin{gather*}
    f(0)=1\\
    f'(0)=k\\
    y-1=kx
\end{gather*}

\problem{4}

\subproblem{a}
\[
L(x)=2+\frac{1}{4}(x-1)
\]

\subproblem{b}
\begin{gather*}
    f'(x)=\frac{1}{2\sqrt{x+3}}\\
    f''(x)=-\frac{1}{4\sqrt{(x+3)^3}}
\end{gather*}
$f(x)$ is concave down on its domain. Hence, both approximations are overestimates.

\problem{5}

\subproblem{a}
\begin{gather*}
    f(x)=x^\frac{2}{3}\\
    f'(x)=\frac{2}{3\cbrt{x}}\\
    L(x)=\frac{1}{3}(x-8)+4\\
    L(8.06)=\boxed{4.02}
\end{gather*}

\problem{6}

\subproblem{b}
\begin{gather*}
    r=30\pm 0.1\\
    V=r^3\\
    dV=3r^2\, dr=3(30)^20.1=\boxed{270}\\
    \frac{dV}{V}=\frac{270}{30^3}=\boxed{\frac{1}{100}}\\
    A=6r^2\\
    dA=12r\, dr=12(30)0.1=\boxed{36}\\
    \frac{dA}{A}=\frac{36}{6\cdot 30^2}=\boxed{\frac{1}{150}}
\end{gather*}

\subproblem{c}
\begin{gather*}
    Q=kx^n\\
    dQ=knx^{n-1}\, dx\\
    \boxed{\frac{dQ}{Q}=n\frac{dx}{x}}
\end{gather*}
The relative error is for a quantity proportional to the $n$-th power of another variable is equal to $n$ times the relative error for that variable.

\problem{7}

\subproblem{a}
\begin{gather*}
    L(x)=2(x-1)+5=2x+3\\
    f(0.9)\approx \boxed{4.8}\\
    f(1.1)\approx \boxed{5.2}
\end{gather*}

\subproblem{b}
As $f'(x)$ is decreasing on its interval, $f''(x)<0$. Hence, the approximations are overestimates. \vem

\subproblem{c}
$f(1)-f(0.9)$ is larger, as the area under $f'(x)$ for $(0.9,1)$ is greater than the area for $(1,1.1)$.\vem

\subproblem{d}
\begin{gather*}
    L(3)=9\\
    f(3)\approx 8
\end{gather*}
I would trust the trapezoidal estimate more, as the slope changes significantly from 1 to 3 and the linearization method is simply not precise.

\problem{8}

\subproblem{a}
\begin{gather*}
    h=50\tan \frac{\pi}{3}=\boxed{50\sqrt{3}}\\
    h=50\tan \theta\\
    dh = 50 \sec^2 \theta \, d\theta \\
    dh=50\cdot 4\cdot \frac{\pi}{360}= \boxed{\frac{5\pi}{9}}\\
    \frac{dh}{h}=\frac{5\pi}{9} \cdot \frac{1}{50\sqrt{3}}=\boxed{\frac{\pi}{90\sqrt{3}}}
\end{gather*}

\subproblem{b}
\begin{gather*}
    h=d\tan \theta \\
    dh = d \sec^2 \theta \, d\theta \\
    \frac{dh}{h}=\frac{\sec^2\theta}{\tan \theta }=\frac{d\theta }{\sin \theta \cos \theta } =\boxed{\frac{2\,d\theta }{\sin 2\theta }}
\end{gather*}

\subproblem{c}
\begin{gather*}
    \sin 2\theta =1\implies \theta =\boxed{\frac{\pi}{4}}\\
    d=h\\
    |d\theta |=\frac{\pi}{360} \implies E=\boxed{\frac{\pi}{180}}
\end{gather*}

\subproblem{d}
When standing either very far away or standing very close to the object, any slight difference in $\theta $ may lead to huge differences in the calculated value.

\problem{9}

\subproblem{a}
\begin{gather*}
    (1+x)^k\approx 1+kx\\
    k=\frac{1}{3}\implies (1-x)^\frac{1}{3}\approx 1-\frac{1}{3}x\\
    k=-4\implies (1+2x)^{-4}\approx 1-8x
\end{gather*}

\subproblem{b}
\begin{gather*}
    |(1-x)^{\frac{1}{3}} - (1 - \tfrac{1}{3}x)| \le 0.1\\
    \iff -1.20 \le x \le 0.71\\
    |(1+2x)^{-4} - (1 - 8x)| \le 0.1 \\
    \iff -0.05 \le x \le 0.06
\end{gather*}
The magnitude of the second function's second derivative is much larger than the first function's near $a=0$.\vem

\subproblem{c}
\begin{gather*}
    K=mc^2\paren{\frac{1}{\sqrt{1-v^2/c^2}}-1}= mc^2\paren{1+\frac{v^2}{2c^2}-1}= \boxed{\frac{1}{2}mv^2}
\end{gather*}
It is an understatement, as the second derivative of the function $f(x)=\frac{1}{\sqrt{1-x}}$ is greater than zero.\vem

\subproblem{d}
\begin{gather*}
    \ln (1)=0\\
    \ddx \ln (1)= 1\\
    y=x\\
    T=\frac{\ln 2}{\ln (1+R/100)} = \frac{100\ln 2}{R}\approx \boxed{\frac{69.3}{R}}
\end{gather*}
It is an underestimate as the second derivative of the linearization function is greater than zero near $R=0$.

\problem{10}

\subproblem{a}
\begin{gather*}
    L(x)=f(x_n)+f'(x_n)(x-x_n)\\
    L(x_{n+1})=0 \\
    f(x_n)+f'(x_n)(x_{n+1}-x_n)=0\\
    \boxed{x_{n+1}=x_n-\frac{f(x_n)}{f'(x_n)}}
\end{gather*}

\subproblem{b}
\begin{gather*}
    f(x)=x^2-2\\
    f'(x)=2x\\
    x_{n+1}=x_n-\frac{x_n^2-2}{2x_n}=\frac{1}{2}\paren{x_n+\frac{2}{x_n}}\\
    x_0 =1\\
    x_1=\frac{3}{2}=1.5\\
    x_2=\frac{17}{12}=1.416\\
    x_3=\frac{577}{408}=1.4142
\end{gather*}

\subproblem{c}
\begin{gather*}
    x_{n+1}=\frac{x_n^2+2}{2x_n}\\
    x_{n+1}-\sqrt{2}=\frac{x_n^2+2}{2x_n}-\sqrt{2}=\frac{x_n^2+2\sqrt2 x_n+2}{2x_n}= \frac{(x_n-\sqrt{2})^2}{2x_n}
\end{gather*}
The absolute error is squared now, rendering the number of correct decimal places to also double. \vem

\subproblem{d}
\begin{gather*}
    x_{n+1}=x_n-\frac{x_n^3-2x_n-5}{3x_n^2-2}\\
    x_0=2\\
    x_1=2.1\\
    x_2=2.09456\\
    x_3=2.09455
\end{gather*}

\subproblem{e}
\begin{gather*}
    x_{n+1}=-2x_n\\
    x_0=1\\
    x_1=-2\\
    x_2=4\\
    x_3=-8\\
\end{gather*}
The sequence oscillates as the derivative of the function approaches infinity as $x$ approaches 0.

\problem{11}

\subproblem{a}
\begin{gather*}
    f(\frac{1}{3})=\frac{3}{10}\\
    f(x)=\frac{x}{1+x^2}\\
    f'(x)=\frac{1-x^2}{(1+x^2)^2}\\
    f(\frac{1}{3})=\frac{\frac{8}{9}}{\frac{100}{81}} = \frac{18}{25}\\
    L(x)=\frac{3}{10}+\frac{18}{25}(x-\frac{1}{3})= \frac{18}{25}x + \frac{3}{50}\\
    f(x)\le L(x)\\
    \frac{x}{1+x^2}\le \frac{18}{25}x+\frac{3}{50}\\
    50x \le (1+x^2)(36x+3)\\
    36x^3+3x^2-14x+3 \ge 0\\
    (3x-1)^2(4x+3)\ge 0
\end{gather*}

\subproblem{b}
\begin{gather*}
    \frac{a}{a^2+1}+\frac{b}{b^2+1}+\frac{c}{c^2+1} \le \frac{18}{25}(a+b+c)+\frac{9}{50} =\frac{18}{25}+\frac{9}{50}= \frac{9}{10}\\
    \boxed{\frac{a}{a^2+1}+\frac{b}{b^2+1}+\frac{c}{c^2+1} \le \frac{9}{10}}\\
    \boxed{a=b=c=\frac{1}{3}}
\end{gather*}

\subproblem{c}
\begin{gather*}
    B(p)=3f(p)+f'(p)(1-3p)\\
    B'(p)=3f'(p)+f''(1-3p)-3f'(p)=f''(1-3p)\\
    f''(p)= \frac{d}{dp}  \frac{1-p^2}{(1+p^2)^2} = \frac{2x(x^2-3)} {(1+p^2)^3}\\
    B'(p)=\frac{2x(x^2-3)(1-3p)} {(1+p^2)^3}\\
    B'(p)=0 \implies \boxed{p=\frac{1}{3}}
\end{gather*}

\subproblem{d}
\begin{gather*}
    b+c=3-a\\
    \frac{1}{a^2+b+c}= \frac{1}{a^2-a+3}\\
    f(x)=\frac{1}{x^2-x+3}\\
    f(1)=\frac{1}{3}\\
    f'(1)= -\frac{1}{9}\\
    L(x)= -\frac{1}{9}(x-1)+\frac{1}{3}=\frac{4-x}{9}\\
    f(x)\le L(x), \quad x \in (0,3)\\
    \frac{1}{a^2+b+c} + \frac{1}{b^2+c+a} + \frac{1}{c^2+a+b} \le  \frac{12 - (a + b + c)}{9} = 1\\
    \boxed{\frac{1}{a^2+b+c} + \frac{1}{b^2+c+a} + \frac{1}{c^2+a+b} \le 1}
\end{gather*}














\end{document}

